\bookchapter{Bab 9: Pejabat \& Pembelajaran Bergerak}

\chapterintro{Saya dahulu juga rasa bekerja di telefon hanyalah “bersusah-susah sedikit”, sehingga satu kali perjalanan kerja terlupa bawa komputer, terpaksa mengubah dokumen tender dan menandatangani kontrak menggunakan telefon. Sejujurnya, halangan kerja di telefon selalunya bukan pada telefon itu sendiri, tetapi kerana anda tidak mengenali batas keupayaannya.}

\booksection{Pemprosesan dokumen: Edit \& cipta di telapak tangan}

: Jika syarikat anda menggunakan Office asli, ini adalah pilihan terbaik. Keserasian format dengan versi desktop hampir sempurna, tidak akan berlaku situasi WPS di mana susun atur kompleks sedikit berubah. Integrasi mendalam dengan OneDrive, dokumen yang anda separuh ubah di komputer, bila dibuka di telefon ia adalah versi terkini.

\begin{itemize}
  \item \textbf{Shimo Documents/Tencent Docs}: Kedua-dua ini adalah wakil "kolaborasi dalam talian". Sesuai untuk ramai orang mengedit dokumen yang sama secara serentak, setiap perubahan dipaparkan secara masa nyata. Saya guna Shimo Documents dengan kawan untuk membuat panduan perjalanan, dua orang mengisi maklumat dalam jadual serentak, jauh lebih cekap daripada seorang buat sendiri.
  \item \textbf{Aliran kerja praktikal}: Tabiat saya sekarang ialah—fail penting disimpan di OneDrive, edit harian guna WPS, kerjasama team guna Shimo Documents. Selepas fail naik ke awan, tidak pernah lagi berlaku situasi malu "terlupa bawa USB drive".
  \item \textbf{2. Apl Nota: Daripada catatan ringkas ke sistem pengetahuan}Banyak orang mahu terus buat sistem yang sangat kompleks apabila mula mencatat nota, tapi tak bertahan tiga hari. Cadangan saya ialah ringan dahulu, berat kemudian—biasakan "mencatat" dahulu, baru fikir tentang "pengurusan".
\end{itemize}

\textbf{Catatan ringkas—sistem terbina dalam sudah memadai}: iOS Notes lancar sangat cepat, dan menyokong imbasan dokumen, tulisan tangan, rakaman. Nota pelbagai jenama Android juga tidak kurang hebat, boleh dilancarkan dengan cepat walaupun skrin terkunci. Alat ini sesuai untuk mencatat idea sementara, senarai beli-belah, perkara penting mesyuarat. Kuncinya lancar pantas, bila terfikir boleh terus catat, tak perlu tunggu App dimuatkan.

\booksection{Apl Nota: Daripada catatan ringkas ke sistem pengetahuan}

: Fungsi “Tetingkap Terapung” WeChat sebenarnya alat tersembunyi. Nampak artikel bagus tapi tak sempat baca, klik “Tetingkap Terapung” di penjuru kanan atas, ia akan mengecil menjadi titik bulat kecil di tepi skrin, boleh diklik bila-bila untuk terus baca. Jauh lebih mudah daripada memajukan ke “Pembantu Pemindahan Fail”. Tambahan pula, alat baca-kemudian seperti Cubox atau Pocket boleh mengumpulkan kandungan bagus yang anda lihat di mana sahaja merentas platform.

\begin{itemize}
  \item \textbf{Bina sistem—Notion dan Lark}: Apabila anda rasa nota berselerak mula sukar dicari, boleh guna Notion. Ia sangat fleksibel—anda boleh buat "pangkalan data nota bacaan", automatik disusun mengikut tajuk buku, penulis, markah; atau bina "papan kanban projek", jadikan kemajuan tugas kerja visual. "Jadual Pelbagai Dimensi" Lark juga serupa, dan kelajuan akses dalam negara lebih cepat. Saya guna Notion untuk menguruskan semua bahan penulisan, nota pembelajaran dan kemajuan projek, sudah lebih tiga tahun.
  \item \textbf{Strategi nota praktikal}: Tabiat saya ialah "tiga set"—idea sementara guna nota sistem (lancar pantas), artikel baca-kemudian guna tetingkap terapung WeChat (senang), pengetahuan jangka panjang letak di Notion (sistematik). Jangan terus mahu masukkan semua benda dalam satu alat, ini sebaliknya mudah membuatkan anda berputus asa.
  \item \textbf{3. Imbasan \& OCR: Masukkan dunia kertas ke dalam telefon}Perkara imbasan telefon, jika digunakan dengan baik, benar-benar menjimatkan banyak masalah pergi ke kedai cetak. Lepas kali saya mohon visa perlukan versi elektronik banyak dokumen, semuanya selesai dengan telefon.
\end{itemize}

\textbf{Pengguna iPhone}: Fungsi imbasan terbina dalam iOS Notes sudah sangat bagus. Buka Notes klik ikon kamera pilih “Imbas Dokumen”, halakan ke fail ia akan automatik kenal pasti bingkai dan betulkan perspektif. Selepas ambil boleh laraskan warna—pilih mod “Hitam Putih” hasil paling hampir dengan pengimbas, “Skala Kelabu” sesuai untuk dokumen bergambar.

\booksection{Imbasan \& OCR: Masukkan dunia kertas ke dalam telefon}

: Kamera banyak jenama ada mod “Imbasan Dokumen” terbina dalam. Jika tidak jumpa, muat turun Microsoft Lens di kedai aplikasi (percuma, dan disambungkan dengan OneDrive).

\begin{itemize}
  \item \textbf{Pengenalan teks OCR}: Fail imbasan, jika perlu ekstrak teks untuk diedit, perlukan OCR. Fungsi "Ekstrak Teks" WeChat paling mudah—tekan lama gambar boleh guna, ketepatan pengenalan agak baik. Jika proses banyak, boleh guna App Baimiao atau CamScanner. Untuk pengenalan bahasa Cina, enjin OCR Baidu lebih baik; fungsi gambar-ke-teks WPS juga sangat tepat.
  \item \textbf{Tip imbasan kecil}: Letakkan di tempat dengan cahaya mencukupi untuk ambil gambar, terbaik ada lampu meja dari sisi, elakkan bayangan tangan pada dokumen. Jika fail berlipat, ratakan dengan tangan dahulu sebelum ambil. Fail berbilang muka surat boleh diimbas berturut-turut kemudian digabungkan menjadi satu PDF, tak perlu simpan satu-satu muka surat.
  \item \textbf{4. Cara pilih storan awan: tiada yang terbaik, hanya yang paling sesuai untuk anda}Saya fikir lama tentang pilihan storan awan, akhirnya mendapati ia bergantung kepada "ekosistem utama" anda.
  \item \textbf{Keluarga Apple iCloud}: Penyegerakan lancar sehingga anda tak perasan kewujudannya. Gambar yang anda ambil di telefon, beberapa minit kemudian automatik ada di komputer. Dan perlindungan privasi Apple paling ketat. Tetapi kelemahannya jelas: berkongsi merentas platform dengan pengguna Android atau Windows agak rumit, percuma hanya 5GB, tambah kapasiti perlukan bayaran.
\end{itemize}

\booksection{Cara pilih storan awan: tiada yang terbaik, hanya yang paling sesuai untuk anda}

: Jika anda guna komputer Windows + Office 365, OneDrive adalah pasangan semula jadi. Simpan dokumen dalam Word, automatik disimpan ke awan. Ruang 1TB dengan langganan Office sangat berbaloi. Fungsi sandaran automatik gambar App telefon juga sangat berguna.

\begin{itemize}
  \item \textbf{Kongsi dalam negara Baidu Netdisk}: Sejujurnya, had laju muat turun Baidu Netdisk sentiasa dikritik, tapi ekosistem kongsinya memang tiada tandingan. Banyak bahan pembelajaran dan sumber filem di internet Cina dikongsi melalui Baidu Netdisk. Jika hanya sekali-sekala muat turun fail yang dikongsi orang, tanpa ahli juga boleh guna. Jika anda pengumpul sumber berat, ahli bayar patut dilanggan, masa yang dijimatkan lebih berharga.
  \item \textbf{Alibaba Cloud Drive/Quark Netdisk}: Dua tahun ini pendatang baru. Alibaba Cloud Drive muat naik turun tanpa had laju (buat masa ini), Quark Netdisk terikat mendalam dengan Quark Browser, semasa mencari bahan boleh terus simpan ke storan awan, sangat mudah.
  \item \textbf{Pelan gabungan saya}: iCloud urus gambar peribadi dan fail harian (privasi keutamaan), Baidu Netdisk simpan sumber dan terima perkongsian (ekosistem keutamaan), OneDrive urus dokumen kerja (integrasi Office keutamaan). Tiga storan awan berfungsi masing-masing, tidak mengganggu antara satu sama lain.
  \item \textbf{5. Senario pejabat bergerak yang sebenar}Cerita pengalaman sebenar saya sendiri. Sekali pergi perjalanan kerja ke luar bandar, pelanggan tiba-tiba mahu ubah syarat kontrak. Aliran kerja saya adalah seperti ini:
\end{itemize}

\textbf{Di hotel menerima cadangan ubahan dari pelanggan melalui WeChat (mesej teks + satu gambar anotasi tulisan tangan).}Buka OneDrive cari kontrak asal, guna WPS buka edit, rujuk mesej WeChat pelanggan ubah syarat satu persatu.

\booksection{Senario pejabat bergerak yang sebenar}

Pelanggan beri pendapat pusingan kedua, saya guna alat pen WPS terus bulatkan kedudukan ubahan pada PDF, dengan anotasi suara jelaskan idea ubahan.

\begin{enumerate}
  \item Selepas sahkan tiada kesilapan, guna fungsi "Tandatangan Elektronik" WPS tandatangan dalam talian, hantar kepada pihak lain untuk tandatangan juga.
  \item Akhir sekali simpan PDF kontrak yang ditandatangani ke Baidu Netdisk dan iCloud satu salinan masing-masing, pada masa yang sama catat dalam log kerja Notion.
  \item Seluruh proses dari mula hingga selesai kira-kira empat puluh minit, sepanjang masa tidak buka komputer. Sejujurnya, tiga tahun lalu saya tidak percaya telefon boleh buat semua ini, tapi sekarang memang boleh. Kuncinya kenali alat-alat ini lebih awal, jangan tunggu bila perlukan baru tergopoh-gopoh muat turun dan daftar.
  \item Cuba ini: Bina aliran kerja pejabat bergerak anda
  \item Pilih Apl dokumen utama anda (WPS atau Microsoft 365), log masuk dan hubungkan dengan storan awan.
  \item Buka Notes atau nota, guna fungsi imbasan ambil satu dokumen, eksport sebagai PDF hantar kepada diri sendiri.
\end{enumerate}

Pilih satu Apl storan awan, tetapkan sandaran automatik gambar telefon (pilih “Sandaran hanya dalam WiFi” jimat data).

\begin{tcbpractice}{Bina satu halaman baharu dalam Notion atau Lark, cuba bina satu templat "Log Kerja" yang mudah.}
\begin{enumerate}
  \item Bab 10: Pembayaran Bergerak \& Perkhidmatan Hidup
  \item Sejujurnya, generasi ibu bapa saya pada mulanya menentang perkara "bayar dengan telefon". "Wang dalam telefon, kalau dicuri orang macam mana?" Kebimbangan ini normal, tapi realitinya: wang tunai dalam poket anda, kebarangkalian hilang jauh lebih tinggi daripada pembayaran telefon dicuri. Pembayaran telefon bukan menyimpan wang dalam telefon, ia hanyalah satu kunci, kunci yang membuka akaun bank anda.
  \item 1. Bagaimana sebenarnya untuk mengikat kad bank?
  \item Sama ada anda guna Alipay atau WeChat, kali pertama mengikat kad perlu lalui langkah ini:
\end{enumerate}

\end{tcbpractice}
